\begin{figure}[htbp]
\centering
\begin{tikzpicture}[x=1cm,y=1cm,font=\small,
 flow/.style={-{Stealth[length=1.5mm]},draw=accent,line width=.65pt}]
\definecolor{ba}{HTML}{8CBCD6}
\definecolor{bb}{HTML}{8FC6B2}
\definecolor{bc}{HTML}{EDBF82}
\foreach \s/\title in {0/{(a) 逐 query 计算},4.9/{(b) 相邻 query 分组},9.8/{(c) 按 KV 块分组}}{
 \begin{scope}[xshift=\s cm]
 \node[font=\bfseries] at (2.15,4.9) {\title};
 \foreach \j/\name in {0/A,1/B,2/C}{\node[font=\footnotesize] at (1.1+\j*.78,4.37) {\name};}
 \foreach \i in {0,1,2}{
  \node[anchor=east,font=\footnotesize] at (.58,3.85-\i*.57) {$q_{\i}$};
  \foreach \j in {0,1,2}{\draw[fill=waste!55,draw=white] (.74+\j*.78,3.61-\i*.57) rectangle ++(.72,.48);}
 }
 \foreach \i/\j/\c in {0/0/ba,0/1/bb,1/1/bb,1/2/bc,2/0/ba,2/2/bc}{
  \filldraw[fill=\c,draw=white] (.74+\j*.78,3.61-\i*.57) rectangle ++(.72,.48);
 }
 \end{scope}
}
% Owner: one row per task, each owns its complete output.
\foreach \i in {0,1,2}{
 \draw[ink,dashed,rounded corners=1pt] (.66,3.55-\i*.57) rectangle ++(2.42,.59);
 \draw[flow] (3.13,3.85-\i*.57) -- (3.56,3.85-\i*.57);
 \draw[fill=selected,draw=accent] (3.64,3.61-\i*.57) rectangle ++(.57,.48);
 \node[font=\scriptsize] at (3.925,3.85-\i*.57) {$O_{\i}$};
}
\node[font=\footnotesize] at (2.17,2.04) {输出不用合并};
% Union: one larger task, holes retain the per-query mask.
\draw[ink,dashed,rounded corners=1pt] (5.56,2.38) rectangle (7.98,4.14);
\foreach \i/\j in {0/2,1/0,2/1}{
 \fill[pattern=north east lines,pattern color=muted!65] (5.64+\j*.78,3.61-\i*.57) rectangle ++(.72,.48);
}
\foreach \i in {0,1,2}{
 \draw[flow] (8.04,3.85-\i*.57) -- (8.46,3.85-\i*.57);
 \draw[fill=selected,draw=accent] (8.54,3.61-\i*.57) rectangle ++(.57,.48);
 \node[font=\scriptsize] at (8.825,3.85-\i*.57) {$O_{\i}$};
}
\node[font=\footnotesize] at (7.05,2.04) {共用 A/B/C；斜线格是无效计算};
% KV owner: split by columns. Highlight A and gather non-adjacent Q rows.
\foreach \j in {0,1,2}{\draw[ink,dashed,rounded corners=1pt] (10.49+\j*.78,2.4) rectangle ++(.83,1.74);}
\draw[accent,line width=1pt] (10.49,2.4) rectangle ++(.83,1.74);
\draw[flow] (10.9,2.35) -- (10.9,1.75);
\foreach \i/\q in {0/0,1/2}{
 \draw[fill=selected,draw=accent] (10.32,1.45-\i*.28) rectangle ++(1.15,.25);
 \node[font=\scriptsize] at (10.895,1.575-\i*.28) {$Q_{\q}$};
 \draw[fill=ba,draw=accent] (12.44,1.45-\i*.28) rectangle ++(1.32,.25);
 \node[font=\scriptsize] at (13.10,1.575-\i*.28) {$R_{\q,A}$};
}
\draw[flow] (11.54,1.43) -- (12.36,1.43);
\node[font=\scriptsize] at (11.95,1.82) {$K_A,V_A$};
\node[font=\footnotesize] at (12.08,.77) {收集选中这块的 query};
\node[font=\footnotesize] at (12.08,.36) {局部结果最后合并};
% Swap is an orthogonal local transformation, not a fourth owner.
\draw[gridline] (0,1.64) -- (9.2,1.64);
\node[anchor=west,font=\footnotesize\bfseries] at (.1,1.23) {Swap AB};
\foreach \r in {0,1}{\foreach \c in {0,1,2,3}{
 \draw[fill=selected,draw=white] (2.2+\c*.34,.53+\r*.25) rectangle ++(.34,.25);
 \draw[fill=selected,draw=white] (4.5+\r*.25,.3+\c*.25) rectangle ++(.25,.25);
}}
\draw[flow] (3.72,.78) -- (4.32,.78);
\node[font=\footnotesize,align=left,anchor=west] at (5.45,.81) {只改 CTA 内部\\的矩阵方向};
\draw[ink,dashed] (.1,-.25) rectangle ++(.4,.23);
\node[anchor=west,font=\footnotesize,text=muted] at (.65,-.135) {虚线：一个 CTA 的范围；彩色：被选中的块；$R$：局部结果};
\end{tikzpicture}
\caption{同一组选块的三种任务划分}
\label{fig:kernel-map}
\end{figure}
